\documentclass[10pt,a4paper]{amsart}
\usepackage[margin=19mm]{geometry}
\usepackage{amsmath,amssymb,mathtools}
\usepackage[scr=rsfs,cal=euler]{mathalfa}
\usepackage{xfrac}
\usepackage[unicode,hidelinks,bookmarks=false]{hyperref}
\newcommand{\ie}{i.e.\ }
\newcommand{\cf}{cf.\ }
\newcommand{\resp}{respectively\ }
\newcommand{\Subsection}{\subsection}
\providecommand{\nobreakdash}{\nobreak\hbox{-}}
\setlength{\emergencystretch}{2em}
\multlinegap0pt
\usepackage{amssymb}
\usepackage{csquotes}
\usepackage{xcolor}
\newtheorem{proposition}{Proposition}
\newtheorem{theorem}{Theorem}
\newtheorem{lemma}{Lemma}
\def\E{\mathbb E}
\def\F{\mathbb F}
\def\P{\mathbb P}
\def\Z{\mathbb Z}
\newcommand{\e}{\varepsilon}
\title{A random polynomial with multiplicative coefficients is almost surely irreducible}
\author{P\'eter P. Varj\'u, Max Wenqiang Xu}
\date{Source: arXiv:2511.04240v2 (CC BY 4.0)}
\begin{document}
\maketitle

\noindent\emph{Source and licence.} A random polynomial with multiplicative coefficients is almost surely irreducible. Version arXiv:2511.04240v2 (CC BY 4.0), distributed by arXiv under the Creative Commons Attribution 4.0 International licence, http://creativecommons.org/licenses/by/4.0/, as recorded in the arXiv licence metadata for this exact version and shown on the abstract page https://arxiv.org/abs/2511.04240v2. Full licence notice: \texttt{licenses/arxiv-2511.04240-CC-BY-4.0.txt}.\par\smallskip

\noindent\textit{Editorial scope.} Counted unit: the article's proposition pr:number-roots-random-poly (source0.tex lines 822--831). The article's complete original proof of it is delivered with it (source0.tex lines 880--967, one \texttt{proof} environment of 88 lines, which the article places away from the statement). No mathematics is written, repaired or re-derived: the statement and the proof are the article's own lines, in the article's own notation, and no retained line is substituted or re-flowed.  Retained uncounted as the source-local context the proof needs: lines 475--494; lines 772--785; lines 855--861; lines 866--870.  Nothing was omitted from any retained span, so the package carries no omission record.  No retained line contains a citation, so the wrapper carries no bibliography and no bibliography entry.  No retained line contains a figure environment, an image-inclusion command or drawing code, so no image is delivered and the package declares no asset dependency.  Editorial formatting only, in addition: an AMS \texttt{amsart} preamble that restates the article's own declarations and macros used by the retained lines, namely the environments \texttt{proposition}, \texttt{theorem}, \texttt{lemma} on separate counters, together with the standard packages those lines need.  The retained environments print in this document as Proposition 1, Theorem 1, Lemma 1 in order of appearance, whereas the article prints the same items with the numbering of its own full document.  The complete native source of the article as distributed in the arXiv e-print for this exact version is bundled unchanged as \texttt{sources/188571/source0.tex}.
\begin{proposition}\label{pr:unexceptional}
Let $l,d\in\Z_{\ge 3}$ and let $q$ be a prime that is suitably large in
terms of $l$.
Let $K\in\Z_{\ge1}$.
Let $a\in\F_q$ be an element such that $a^k$
is not $l$-exceptional for any $k=1,\ldots,K$.
Let $I\subset [0,\ldots,d]$ be a set that contains $q^{5/(l+1)}$ pairwise disjoint
arithmetic progressions of length $3l^3$ with common difference at most $K$.
Let $X_i$ be independent uniform $\pm1$ valued random variables, and
let
\[
Y=\sum_{i\in I}X_i(a^i,ia^{i-1})
\]
be a random element of $\F_q^2$.
Then
\[
\Big|\P[Y=x]-\frac{1}{q^2}\Big|< q^{-10}
\]
for all $x\in\F_q^2$.
\end{proposition}

\begin{proposition}\label{pr:exceptional}
Let $q$ be an odd prime, and let $a\in\F_q^\times$.
Let $I$ be a set of positive integers with $|I|\le q$.
Let $X_i$ be independent uniform $\pm1$ valued random variables for $i\in I$, and let
\[
Y=\sum_{i\in I} X_i a^i
\]
be a random element of $\F_q$.
Then
\[
\P[Y=x]\le 129 |I|^{-1/2}.
\]
for all $x\in \F_q$.
\end{proposition}

\begin{theorem}\label{th:Green-Tao}
For every $L\in\Z_{>0}$, there is a constant $C=C(L)$ such that the following
holds for all $d\in\Z_{>0}$ that is sufficiently large in terms of $L$.
Let $A$ be a subset of the primes in $[1,d]$ with $|A|>d/10\log d$.
Then $A$ contains an arithmetic progression of length $L$ with common difference
less than $C(\log d)^C$.
\end{theorem}

\begin{lemma}\label{lm:number-exceptional}
For every $l$, there is a constant $C=C(l)$ such that the number of
$(l,q)$-exceptional polynomials is at most
$Cq^{1/(l+1)}$.
\end{lemma}

\begin{proposition}\label{pr:number-roots-random-poly}
Fix $l\in\Z_{>0}$ and $\e\in(0,1/10)$.
Let $q$ be an odd prime and $d\in\Z_{>0}$ that are both sufficiently
large in terms of $l$ and such that $l\ge \e^{-1}\log q/\log d$ and $q>d$.
Then
\begin{align*}
|\E[B_R(q)]-1|&<d^{-1/2+\e},\\
\E[|\{a\in\F_q: \text{$a$ is a double root of $R$}\}|]&<d^{-1/2+\e}.
\end{align*}
\end{proposition}

\begin{proof}[Proof of Proposition \ref{pr:number-roots-random-poly}]
Let $K=C_1(\log d)^{C_1}$,
where $C_1$ is the constant $C$
in Theorem \ref{th:Green-Tao}
applied with $L=3l^3$.
We first consider the probability that some $a\in\F_q$
is a root or a double root of $R$
under the condition that $a^k$ is not $l$-exceptional for any $1\le k\le K$.
To this end, we will apply Proposition \ref{pr:unexceptional} with the set of
primes in $[d/2,d]$ in the role of $I$.


We first show that $I$ contains the required number of arithmetic progressions.
We apply Theorem \ref{th:Green-Tao} repeatedly to find arithmetic progressions
of length $3l^3$ with common difference at most $K$.
We start this with all primes in $[d/2,d]$ in the role of $A$ at first, then
we remove from $A$ all arithmetic progressions that we find to apply
Theorem \ref{th:Green-Tao} again for this reduced set.

This process can be run more than $d/(10\log(d) l^3)$ times before the number
of elements in $A$ falls below the threshold in the theorem, so
we find at least this many arithmetic progressions.
Since $l\ge 5\log q/\log d$,
\[
q^{5/(l+1)}\le \frac{d}{10\log(d) l^3}, 
\]
and we have enough arithmetic progressions to apply Proposition~\ref{pr:unexceptional}.

Next we write
\[
(R(a),R'(a))= Y+Z,
\]
where
\[
Y=\sum_{i\in I}X_{i}(a^{i-1},(i-1)a^{i-2})
\]
and
\[
Z=\sum_{i\in [1,d+1]\backslash I}X_{i}(a^{i-1},(i-1)a^{i-2}).
\]
We observe that $Y$ and $Z$ are independent. So, conditioning on the value of $Z$, we can write
\[
\Big|\P[R(a)=x_1,R'(a)=x_2]-\frac{1}{q^2}\Big|
\le \max_{z\in\F_q^2}\Big|\P[Y=(x_1,x_2)-z]-\frac{1}{q^2}\Big|
< q^{-10}
\]
for all $x_1,x_2\in\F_q$, where we applied Proposition~\ref{pr:unexceptional}.
We conclude
\begin{align}
\Big|\P[R(a)=0]-\frac{1}{q}\Big|&\le q^{-9}\label{eq:unexceptional1}\\
\P[\text{$a$ is a double root of $R$}]\le \frac{2}{q^2}.\label{eq:unexceptional2}
\end{align}

Now we prove a bound that is valid for all $a\in\F_q^{\times}$.
We write $R(a)=Y_1+Z_1$, where $Y_1$ and $Z_1$ are the first components
of the vectors $Y$ and $Z$ defined above.
Conditioning on the value of $Z_1$ and then using
Proposition \ref{pr:exceptional} we can write
\[
\P[R(a)=0]\le \max_{z\in\F_q}\P[Y_1=-z]\le 129 |I|^{-1/2}\le 1000 (d/\log d) ^{-1/2}.
\]
Note that $R(0)=1$ almost surely, so the probability that $0$ is a root is $0$, so
the above bound is valid even for $a=0$.

The number of elements $a\in\F_q$ for which the bounds \eqref{eq:unexceptional1}
and \eqref{eq:unexceptional2} do not apply is at most
$K^2lC_2q^{1/(l+1)}$, where $C_2$ is the constant
$C$ in Lemma \ref{lm:number-exceptional}.
We see that the expected number of exceptional roots of $R$ is less than
\[
1000C_1^2C_2l(\log d)^{2C_1+1/2}q^{1/(l+1)} d^{-1/2}<d^{-1/2+\e}/2
\]
because $l\ge \e^{-1}\log q/\log d$.

Summing up \eqref{eq:unexceptional2} for $a\in\F_q$ that are not exceptional
and combining with the above bound, we get
\[
\E[|\{a\in\F_q: \text{$a$ is a double root of $R$}\}|]\le \frac{2}{q}+d^{-1/2+\e}/2
< d^{-1/2+\e}.
\]

Summing up \eqref{eq:unexceptional1} and combining with the bound for exceptional
roots, we get
\[
\Big|\E[|\{a\in\F_q: R(a)=0\}|]-\frac{q}{q}\Big|
\le \frac{K^2lC_2q^{1/(l+1)}}{q}+q^{-8}+ d^{-1/2+\e}/2<d^{-1/2+\e}.
\]
\end{proof}

\end{document}
